% Executive Summary
\section*{Executive Summary}
\addcontentsline{toc}{section}{Executive Summary}
This report presents a definitive technical review and resolution of 19 critical inconsistencies, logical gaps, and implementation ambiguities identified across two foundational documents: \texttt{PINN\_architecture\_plan.tex} and \texttt{CTL-instruction-final.tex}. These documents collectively outline a sophisticated framework for training a Physics-Informed Neural Network (PINN) on the chaotic three-body problem using a Curriculum-Transfer Learning (CTL) methodology. The identified issues, if left unaddressed, would preclude a successful, robust, and reproducible implementation of the proposed system.

The resolutions provided herein are synthesized into a single, coherent, and actionable blueprint. The analysis is structured around five core themes, each representing a critical stage of the implementation workflow. The key resolutions are summarized as follows:
\begin{enumerate}
    \item \textbf{Foundational Physical Metrics:} The report establishes robust, computable definitions and provides reference implementations for the Lyapunov time ($\tau_L$) and the normalized angular momentum ($L/L_{\max}$), which are essential for the quantitative stratification of training data into curriculum stages. A scalable, parallelized pipeline for computing global normalization constants across the entire dataset is also detailed.
    \item \textbf{Data Scaling and Normalization:} A unified protocol for data transformation is established. This includes standardizing spatiotemporal normalization to a consistent [-1, 1] range, resolving a critical five-order-of-magnitude discrepancy in the gravitational softening parameter, and providing a definitive, consolidated schema for the construction of the 14-dimensional network input vector.
    \item \textbf{Network Architecture and Initialization:} Ambiguities con\-cerning the dynamic scheduling of the SIREN network's frequency parameter ($\omega_0$) are resolved. A novel, curriculum-aware discrete schedule is proposed, directly linking this critical architectural hyperparameter to the complexity of the data at each training stage.
    \item \textbf{Loss Function Framework:} A rigorous and unambiguous specification for the multi-component loss function is provided. This includes enforcing dimensional consistency across all terms, defining a clear state management protocol for the Augmented Lagrangian Method (ALM) parameters, providing a reference implementation for computing kinematics via automatic differentiation, and specifying a physically appropriate parameter for the Huber loss function.
    \item \textbf{Training and Data Sampling Strategy:} Conflicting instructions for data batching and collocation point sampling are reconciled into a single, efficient, and standard methodology. The report formalizes the adoption of the superior curriculum-aware batching protocol and integrates a physically motivated, stage-dependent schedule for adaptive gradient clipping.
\end{enumerate}
The solutions presented are rigorously grounded in established principles from celestial mechanics, numerical analysis, and deep learning. The resulting blueprint ensures that the final CTL-PINN framework is not only physically sound and mathematically consistent but also computationally robust and practically implementable.

\clearpage
